% Part: first-order-logic
% Chapter: tableaux
% Section: soundness-identity

\documentclass[../../../include/open-logic-section]{subfiles}

\begin{document}

\olfileid{fol}{tab}{sid}

\olsection{\usetoken{S}{identity}-সহ বিশুদ্ধতা}

\begin{prop}
অভিন্নতার বিধি-সহ !!^{tableau}s বিশুদ্ধ: কোনো বদ্ধ !!{tableau}
পরিতৃপ্তিযোগ্য নয়।
\end{prop}

\begin{proof}
আগের মতো শুধু দেখাতে হবে, কোনো !!a{tableau}-তে একটি পরিতৃপ্তিযোগ্য
শাখা থাকলে $\eq$-এর কোনো একটি বিধি প্রয়োগের ফলে উৎপন্ন শাখাটিও
পরিতৃপ্তিযোগ্য। শাখার চিহ্নিত !!{formula}s-এর সমষ্টি~$\Gamma$ এবং
~$\Gamma$-কে পরিতৃপ্তকারী কোনো !!a{structure}~$\Struct{M}$ ধরি।

ধরা যাক~$\eq$ ব্যবহার করে, অর্থাৎ চিহ্নিত
!!{formula}~$\sFmla{\True}{\eq[t][t]}$ যোগ করে শাখাটি প্রসারিত
হয়েছে। স্পষ্টতই $\Sat{M}{\eq[t][t]}$; তাই $\Struct{M}$
$\Gamma \cup \{\sFmla{\True}{\eq[t][t]}\}$-কেও পরিতৃপ্ত করে।

$\TRule{\True}{\eq}$ ব্যবহার করে শাখাটি প্রসারিত হলে চিহ্নিত
!!{formula}~$\sFmla{\True}{!A(t_2)}$ যোগ করি; আর $\Gamma$-তে
$\sFmla{\True}{\eq[t_1][t_2]}$ ও $\sFmla{\True}{!A(t_1)}$ উভয়ই
আছে। তাই $\Sat{M}{\eq[t_1][t_2]}$ এবং
$\Sat{M}{!A(t_1)}$। এমন একটি চলরাশি-আরোপ $s$ নিই যেন
$s(x) = \Value{t_1}{M}$। \olref[syn][ass]{prop:sentence-sat-true}
অনুসারে $\Sat{M}{!A(t_1)}[s]$। যেহেতু $\varAssign{s}{s}{x}$,
\olref[syn][ext]{prop:ext-formulas} অনুসারে
$\Sat{M}{!A(x)}[s]$। আবার $\Sat{M}{\eq[t_1][t_2]}$ বলে
$\Value{t_1}{M} = \Value{t_2}{M}$; অতএব
$s(x) = \Value{t_2}{M}$। \olref[syn][ext]{prop:ext-formulas} আবার
প্রয়োগ করে $\Sat{M}{!A(t_2)}[s]$-ও পাই।
\olref[syn][ass]{prop:sentence-sat-true} অনুসারে
$\Sat{M}{!A(t_2)}$। $\TRule{\False}{\eq}$-এর ক্ষেত্রটিও একইভাবে
প্রমাণিত হয়।
% BN-SRC-057 TARGET-CORRECTION-BEGIN
\par\smallskip\noindent\textnormal{\small\textbf{উৎস-সংশোধন (BN-SRC-057):} হিমায়িত ইংরেজি $\TRule{\True}{\eq}$-এর বিশুদ্ধতা-প্রমাণে যোগ করা সূত্রটিকে $\sFmla{S}{!A(t_2)}$ বলা হলেও পূর্বধারণা এবং প্রমাণটি সত্য-চিহ্নিত ক্ষেত্রটিই বিবেচনা করে। বিধির সংজ্ঞা অনুসারে বাংলা লক্ষ্যে $\sFmla{\True}{!A(t_2)}$ পুনরুদ্ধার করা হয়েছে; মিথ্যা-চিহ্নিত বিধিটি পরের বাক্যে আলাদাভাবে উল্লেখিত।} % BN-SRC-057 TARGET-CORRECTION-END
\end{proof}

\end{document}
